\documentclass[11pt]{article}
\usepackage{mathblatt}
\begin{document}
\weit
\typeout{PROBE-BEGIN brueche-dezimalzahlen-e1-k1-s6-v15}
\begin{aufgabe}{\detokenize{brueche-dezimalzahlen-e1-k1-s6-v15}}
\begin{teile}
\teil Probe

\raute{4}{4}
\end{teile}
\end{aufgabe}
\typeout{PROBE-END brueche-dezimalzahlen-e1-k1-s6-v15}
\typeout{PROBE-BEGIN flaechen-e4-k1-s0-v3}
\begin{aufgabe}{\detokenize{flaechen-e4-k1-s0-v3}}
\begin{teile}
\teil Probe

\raute[achsen=diagonalen,diagonalen]{4}{2.5}
\end{teile}
\end{aufgabe}
\typeout{PROBE-END flaechen-e4-k1-s0-v3}
\typeout{PROBE-BEGIN pythagoras-e3-k1-s0-v3}
\begin{aufgabe}{\detokenize{pythagoras-e3-k1-s0-v3}}
\begin{teile}
\teil Probe

\raute[diagonalen]{4}{2.5}
\end{teile}
\end{aufgabe}
\typeout{PROBE-END pythagoras-e3-k1-s0-v3}
\typeout{PROBE-BEGIN symmetrie-abbildungen-e2-k1-s0-v2}
\begin{aufgabe}{\detokenize{symmetrie-abbildungen-e2-k1-s0-v2}}
\begin{teile}
\teil Probe

\raute{4}{2.5}
\end{teile}
\end{aufgabe}
\typeout{PROBE-END symmetrie-abbildungen-e2-k1-s0-v2}
\typeout{PROBE-BEGIN symmetrie-abbildungen-e2-k2-s0-v2}
\begin{aufgabe}{\detokenize{symmetrie-abbildungen-e2-k2-s0-v2}}
\begin{teile}
\teil Probe

\raute[achsen=diagonalen]{5}{3}
\end{teile}
\end{aufgabe}
\typeout{PROBE-END symmetrie-abbildungen-e2-k2-s0-v2}
\end{document}
